\section{Forward Problem}

The first stage of this work focuses on the forward formulation of the one-dimensional heat diffusion equation. In the forward problem, the thermal diffusivity coefficient $\alpha$ is assumed to be known, and the objective is to reconstruct the temperature field $T(x,t)$ satisfying the governing partial differential equation together with the prescribed initial and boundary conditions.

In the context of Physics-Informed Neural Networks (PINNs), the solution is approximated by a neural network whose training is guided by the physics of the problem. More specifically, the network is trained such that its predictions satisfy the heat equation while also respecting the imposed conditions on the boundaries and at the initial time.

The PINN framework allows different training configurations. In a purely physics-driven setting, the model can be trained only from the residual of the governing equation and the associated constraints, without requiring any observational data. Alternatively, synthetic or experimental data points may also be incorporated into the loss function in order to guide the training process and improve convergence.

The loss function typically combines several contributions:

\begin{itemize}
    \item a physics loss enforcing the partial differential equation,
    \item an initial condition loss,
    \item a boundary condition loss,
    \item and optionally a data loss term when observations are available.
\end{itemize}

For the one-dimensional heat equation, the physics residual is defined as:

\begin{equation}
R(x,t)
=
\frac{\partial T}{\partial t}
-
\alpha
\frac{\partial^2 T}{\partial x^2}
\end{equation}

The PINN is therefore optimized such that the residual approaches zero over the spatio-temporal domain.

Although the forward problem is an important benchmark for validating the implementation and understanding the training dynamics of PINNs, it should be noted that PINNs generally do not outperform classical numerical solvers for this type of problem. Traditional methods such as finite difference, finite volume, or finite element methods are typically faster, more accurate, and computationally more efficient for standard forward simulations of diffusion equations.

Consequently, the main interest of studying the forward problem within the PINN framework is not necessarily to replace classical solvers, but rather to better understand how neural networks learn physical constraints and how the optimization process behaves under different training conditions.

In particular, the forward problem provides a controlled environment for investigating:
\begin{itemize}
    \item the balance between physics and data losses,
    \item convergence properties of the optimization process,
    \item the influence of collocation point distributions,
    \item the effect of network architecture and hyperparameters,
    \item and the role of automatic differentiation in enforcing physical laws.
\end{itemize}

This analysis constitutes an essential preliminary step before addressing the inverse problem, which is generally more challenging and represents one of the main motivations for using PINNs in scientific machine learning applications.

% =========================================================
\subsection{Normalization Strategies}
% =========================================================

An important aspect investigated during this work concerns the normalization of the physical variables and governing equations.

Normalization is commonly employed in Physics-Informed Neural Networks in order to:
\begin{itemize}
    \item improve numerical stability,
    \item reduce scale disparities between variables,
    \item accelerate optimization,
    \item and improve gradient propagation during training.
\end{itemize}

However, for inverse problems involving unknown physical parameters, normalization may also unintentionally modify the sensitivity of the optimization process and introduce hidden dependencies between variables and trainable parameters.

For this reason, several normalization strategies were explored and compared throughout this work.

% =========================================================
\subsubsection{Fully Non-Dimensionalized Formulation}
% =========================================================

The first approach consists of fully non-dimensionalizing both space and time variables using characteristic scales.

Defining:

\begin{equation}
x^*=\frac{x}{L}
\end{equation}

and

\begin{equation}
t^*=\frac{\alpha t}{L^2}
\end{equation}

the heat equation becomes:

\begin{equation}
\frac{\partial T}{\partial t^*}
=
\frac{\partial^2 T}{\partial x^{*2}}
\end{equation}

In this formulation, the diffusion coefficient $\alpha$ disappears from the PDE.

Although this normalization is mathematically elegant and commonly used in classical physics, it becomes problematic for inverse PINN problems because the parameter to identify no longer explicitly appears inside the governing equation.

As a consequence:
\begin{itemize}
    \item the optimization may lose direct sensitivity to $\alpha$,
    \item parameter identifiability may deteriorate,
    \item and the inverse problem becomes artificially modified.
\end{itemize}

This formulation was therefore found to be unsuitable for the present inverse identification problem.

% =========================================================
\subsubsection{Partially Normalized Formulation}
% =========================================================

A second strategy consists of normalizing only the spatial and temporal coordinates independently from $\alpha$.

For example:

\begin{equation}
x^*=\frac{x}{L}
\end{equation}

\begin{equation}
t^*=\frac{t}{t_{max}}
\end{equation}

In this case, the governing equation becomes:

\begin{equation}
\frac{\partial T}{\partial t^*}
=
\alpha
\frac{t_{max}}{L^2}
\frac{\partial^2 T}{\partial x^{*2}}
\end{equation}

Unlike the fully non-dimensionalized formulation, $\alpha$ remains explicitly present in the PDE residual.

This approach improves:
\begin{itemize}
    \item numerical conditioning,
    \item optimization stability,
    \item and input scaling consistency,
\end{itemize}

while preserving the inverse identification structure of the problem.

However, the scaling coefficients introduced in the PDE residual still modify the relative magnitude of the different loss terms and may influence optimization dynamics.

% =========================================================
\subsubsection{Non-Normalized Physical Formulation}
% =========================================================

A third approach consists of directly solving the PDE using the original physical variables without normalization.

The governing equation remains:

\begin{equation}
\frac{\partial T}{\partial t}
=
\alpha
\frac{\partial^2 T}{\partial x^2}
\end{equation}

Although this formulation may initially appear less numerically convenient, it presents several important advantages for inverse parameter estimation.

Most importantly:
\begin{itemize}
    \item the physical meaning of the PDE is fully preserved,
    \item $\alpha$ remains directly coupled to the diffusion dynamics,
    \item and no artificial scaling modifies the identifiability structure of the inverse problem.
\end{itemize}

During the experiments performed in this work, the non-normalized formulation often provided more interpretable and physically consistent optimization behavior for inverse parameter reconstruction.

In particular, it avoided introducing hidden dependencies between:
\begin{itemize}
    \item normalization scales,
    \item PDE residual magnitudes,
    \item and parameter sensitivities.
\end{itemize}

% =========================================================
\subsubsection{Influence of Normalization on Inverse Problems}
% =========================================================

These experiments highlighted an important observation:

\begin{quote}
Normalization strategies that are beneficial for forward problems are not necessarily appropriate for inverse PINN problems.
\end{quote}

For forward simulations, removing physical scales may improve optimization efficiency without fundamentally altering the problem.

However, for inverse problems:
\begin{itemize}
    \item the unknown parameters themselves define part of the physical scales,
    \item and normalization may unintentionally suppress informative parameter dependencies.
\end{itemize}

This becomes particularly important when studying:
\begin{itemize}
    \item parameter identifiability,
    \item sensitivity analysis,
    \item noisy inverse reconstruction,
    \item and adaptive weighting strategies.
\end{itemize}

The normalization strategy therefore constitutes an important design choice in inverse PINN implementations and should be selected carefully depending on the objectives of the study.